% =====================================================================
% Mirante dos Dados — Working Paper #3
% Tratamento Cirúrgico da Incontinência Urinária no SUS (2015–2020)
%
% Autores: Tatieli da Silva (pesquisa original, 2022) + Leonardo Chalhoub
%          (reprodução, extensão e publicação aberta, 2026)
% Compilação: pdflatex articles/uropro-incontinencia-urinaria.tex (2x)
% =====================================================================

\documentclass[12pt,a4paper,oneside]{article}

\usepackage[T1]{fontenc}
\usepackage[utf8]{inputenc}
\usepackage[brazilian]{babel}
\usepackage{lmodern}
\usepackage[a4paper,top=3cm,bottom=2cm,left=3cm,right=2cm]{geometry}
\usepackage{setspace}
\usepackage{indentfirst}
\usepackage{titlesec}
\usepackage{caption}
\usepackage{booktabs}
\usepackage{array}
\usepackage{multirow}
\usepackage{graphicx}
\usepackage[hidelinks]{hyperref}
\usepackage{url}
\usepackage{float}

\graphicspath{{figures-uropro/}}

\onehalfspacing
\setlength{\parindent}{1.25cm}
\setlength{\parskip}{0pt}

\titleformat{\section}{\normalfont\bfseries\large\uppercase}{\thesection}{0.6em}{}
\titleformat{\subsection}{\normalfont\bfseries\normalsize}{\thesubsection}{0.6em}{}
\titlespacing*{\section}{0pt}{18pt}{8pt}
\titlespacing*{\subsection}{0pt}{12pt}{6pt}

\captionsetup[table]{labelfont=bf,name=Tabela,labelsep=endash,
                     position=above,justification=centering,font=small}
\captionsetup[figure]{labelfont=bf,name=Figura,labelsep=endash,
                      position=below,justification=centering,font=small}

\newcommand{\rs}{R\$\,}

\begin{document}

% ---------- CAPA ---------------------------------------------------------
\begin{titlepage}
  \begin{center}
    \vspace*{1.5cm}
    {\small\textsc{Mirante dos Dados \quad\textbar\quad Working Paper n. 3}}\\
    {\small Versão 1.0 \textemdash{} Abril de 2026}

    \vfill

    {\LARGE\bfseries
     TRATAMENTO CIRÚRGICO DA INCONTINÊNCIA URINÁRIA
     NO SISTEMA ÚNICO DE SAÚDE (2015\textendash 2020):
     VOLUMES, DESPESA, PERMANÊNCIA E DISTRIBUIÇÃO
     GEOGRÁFICA POR VIA DE ACESSO}

    \vspace{1cm}
    {\large\itshape
     Surgical treatment of urinary incontinence under the Brazilian
     Unified Health System (2015\textendash 2020): volumes, expenditure,
     length of stay, and geographic distribution by surgical approach}

    \vfill

    {\large\textbf{Tatieli da Silva}\textsuperscript{1} \quad\textbullet\quad
     \textbf{Leonardo Chalhoub}\textsuperscript{2}}\\[6pt]
    {\small\textsuperscript{1}\,Pesquisa original \textemdash{}
     especialização em Enfermagem (2022).}\\[2pt]
    {\small\textsuperscript{2}\,Reprodução, extensão e publicação aberta
     \textemdash{} engenharia de dados.}\\[6pt]
    {\small\href{mailto:leonardochalhoub@gmail.com}{leonardochalhoub@gmail.com}
     \quad\textbullet\quad
     \href{https://github.com/leonardochalhoub}{github.com/leonardochalhoub}}

    \vfill

    {\small Brasil \\ Abril de 2026}
  \end{center}
\end{titlepage}

% ---------- RESUMO -------------------------------------------------------
\section*{Resumo}
\addcontentsline{toc}{section}{Resumo}
\begin{singlespace}\small\noindent
Este artigo analisa empiricamente o tratamento cirúrgico da incontinência
urinária (IU) no Sistema Único de Saúde brasileiro entre 2015 e 2020,
reproduzindo e estendendo a pesquisa original de Tatieli da Silva
(especialização em Enfermagem, 2022). A unidade analítica é a
\textbf{Autorização de Internação Hospitalar (AIH)} aprovada pelo
SIH-SUS para os procedimentos SIGTAP \texttt{0409010499} (via abdominal)
e \texttt{0409070270} (via vaginal). O artigo documenta cinco achados
principais. \textbf{Primeiro}, o volume de procedimentos via vaginal
(30.688 AIH em 2015\textendash 2020) supera em \textbf{6{,}86 vezes} o
da via abdominal (4.472 AIH), proporção mantida em todos os anos da
série e compatível com as recomendações internacionais que privilegiam
abordagens minimamente invasivas. \textbf{Segundo}, a permanência
hospitalar média da via vaginal (2{,}16 dias) é \textbf{37{,}8\%
menor} que a da via abdominal (3{,}47 dias) \textemdash{} achado
consistente com a literatura clínica que documenta menor invasividade
e recuperação mais rápida da abordagem vaginal. \textbf{Terceiro}, o
custo médio por AIH é praticamente equivalente entre as vias
(\rs 442{,}61 ABD vs.\ \rs 424{,}79 VAG, diferença de \textbf{+4{,}2\%}),
apesar da permanência 60\% mais longa na via abdominal \textemdash{}
reflexo da estrutura de remuneração fixa por procedimento da tabela
SIGTAP-SUS, em que a permanência subjacente é apenas parcialmente
refletida nos serviços hospitalares. \textbf{Quarto}, há concentração
geográfica extrema, especialmente na via vaginal: \textbf{São Paulo
concentra 70{,}3\%} das AIH vaginais e 29{,}8\% das abdominais no
período, indicando provável centralização em hospitais-referência.
\textbf{Quinto}, observa-se queda de \textbf{aproximadamente 50\%}
em ambas as vias em 2020, compatível com o adiamento de cirurgias
eletivas durante a pandemia de COVID-19. O gasto público total
acumulado no período aproximou-se de \textbf{\rs 15{,}0 milhões}
(valores nominais), distribuído em 35.160 internações aprovadas. O
trabalho contribui metodologicamente ao integrar a análise ao pipeline
\textit{open-source} \textit{Mirante dos Dados} em arquitetura
medallion, viabilizando atualização mensal automatizada e
reproducibilidade por terceiros.

\vspace{0.4cm}\noindent\textbf{Palavras-chave:}
incontinência urinária; SUS; SIH-SUS; cirurgia uroginecológica;
análise espacial em saúde; políticas de saúde; dados abertos.
\end{singlespace}

% ---------- ABSTRACT -----------------------------------------------------
\section*{Abstract}
\addcontentsline{toc}{section}{Abstract}
\begin{singlespace}\small\noindent
This paper empirically analyzes the surgical treatment of urinary
incontinence (UI) within the Brazilian Unified Health System (SUS)
between 2015 and 2020, reproducing and extending the original research
by Tatieli da Silva (graduate specialization thesis in Nursing, 2022).
The analytical unit is the approved Hospital Admission Authorization
(AIH) for SIGTAP procedures 0409010499 (abdominal route) and
0409070270 (vaginal route). The paper documents five main findings.
\textbf{First}, the vaginal route is performed \textbf{6.86 times more
frequently} than the abdominal one (30,688 vs.\ 4,472 AIH in
2015\textendash 2020), consistent with international urogynecological
guidelines. \textbf{Second}, mean length of stay is \textbf{37.8\%
shorter} for the vaginal route (2.16 vs.\ 3.47 days). \textbf{Third},
average cost per AIH is essentially equivalent across approaches
(BRL 442.61 vs.\ BRL 424.79, +4.2\% for abdominal), despite the
abdominal route having 60\% longer hospital stays. \textbf{Fourth},
geographic concentration is extreme, especially for the vaginal
route: São Paulo accounts for \textbf{70.3\%} of vaginal and 29.8\%
of abdominal national volume. \textbf{Fifth}, both routes drop
\textbf{approximately 50\%} in 2020 due to elective surgery deferral
during the COVID-19 pandemic. Total public expenditure approached
BRL 15.0 million (nominal). The work contributes methodologically by
integrating the analysis into the open-source \textit{Mirante dos
Dados} medallion-architecture pipeline.

\vspace{0.4cm}\noindent\textbf{Keywords:}
urinary incontinence; SUS; hospital information system; urogynecological
surgery; spatial health analysis; health policy; open data.
\end{singlespace}

\newpage

% ---------- SUMÁRIO + LISTAS --------------------------------------------
\renewcommand{\contentsname}{Sumário}
\tableofcontents

\newpage
\renewcommand{\listfigurename}{Lista de Figuras}
\listoffigures

\vspace{1cm}
\renewcommand{\listtablename}{Lista de Tabelas}
\listoftables

\newpage

% =====================================================================
\section{Introdução}\label{sec:intro}
% =====================================================================

A incontinência urinária (IU) é definida pela \textit{International
Continence Society} (ICS) como qualquer perda involuntária de urina
(HAYLEN \textit{et al.}, 2010), condição de elevada prevalência sobretudo
em mulheres adultas e idosas. Estimativas globais situam a prevalência
feminina entre 25\% e 45\%, com impacto substantivo sobre qualidade de
vida, saúde mental e custos diretos e indiretos para os sistemas de
saúde (HU \textit{et al.}, 2003). No Brasil, a despeito da relevância
demográfica e clínica, a IU permanece subdiagnosticada e subtratada,
parcialmente em razão do estigma associado e da escassez de equipes
uroginecológicas em parcela significativa do território nacional.

O Sistema Único de Saúde (SUS) oferece, por meio do Sistema de
Informações Hospitalares (SIH-SUS), cobertura cirúrgica para os casos
refratários ao tratamento conservador. As intervenções cirúrgicas
disponíveis enquadram-se em duas vias técnicas principais: a
\textbf{via abdominal} (cirurgias tipo Burch e variantes), de maior
porte e geralmente reservada a casos complexos ou com indicações
específicas; e a \textbf{via vaginal} (\textit{slings} mediouretrais e
técnicas correlatas), minimamente invasiva e considerada de primeira
linha pela maioria das diretrizes contemporâneas (FORD
\textit{et al.}, 2017; CHAPPLE \textit{et al.}, 2020).

O presente artigo retoma e atualiza a investigação empírica conduzida
por Tatieli da Silva como trabalho de conclusão de especialização em
Enfermagem (TATIELI, 2022). A pesquisa original analisou as AIH
aprovadas no SIH-SUS para os SIGTAPs \texttt{0409010499} e
\texttt{0409070270} no período 2015\textendash 2020, utilizando
agregados pré-computados do TabNet/DATASUS \textemdash{} escolha
metodologicamente conservadora, mas com limitações conhecidas em
termos de granularidade e atualizabilidade. Este trabalho integra a
análise ao pipeline aberto \textit{Mirante dos Dados}, partindo dos
microdados SIH-RD, o que permite \textit{refresh} mensal automatizado
e desagregação arbitrária para investigações subsequentes.

A linha do tempo do contexto clínico-institucional analisado é
sintetizada na Figura \ref{fig:timeline}.

\begin{figure}[H]
\centering
\includegraphics[width=0.95\linewidth]{fig01-timeline-uropro}
\caption{Linha do tempo do contexto clínico-institucional do estudo}
\label{fig:timeline}
\par\noindent{\small Fonte: elaboração própria. SIGTAP padronizado em 2008
(MS); diretriz IUGA/ICS publicada em 2010; \textit{Cochrane Review}
sobre \textit{slings} em 2017.}
\end{figure}

Pretende-se, assim, alcançar três objetivos. \textbf{Primeiro},
confirmar empiricamente os achados centrais da pesquisa original
quanto à distribuição geográfica e ao volume de procedimentos.
\textbf{Segundo}, ampliar o escopo analítico ao quantificar a
permanência hospitalar, o custo médio por AIH e o impacto da
pandemia de COVID-19. \textbf{Terceiro}, oferecer infraestrutura
aberta de dados e visualização que reduza o custo marginal de
pesquisa para profissionais de Enfermagem, gestores hospitalares
e pesquisadores em saúde coletiva.

Este artigo é o terceiro da série de \textit{Working Papers} do
projeto Mirante dos Dados (CHALHOUB, 2026a; CHALHOUB, 2026b). Os três
trabalhos compartilham infraestrutura técnica, dimensões auxiliares
(IBGE/SIDRA, IPCA-BCB) e princípio metodológico \textemdash{} partir
dos microdados oficiais para preservar flexibilidade analítica.
Diferentemente dos dois primeiros, em que a unidade primária é
monetária (\rs pago) e a lógica alocativa é institucional-política
ou socioeconômica, este terceiro estudo tem a internação cirúrgica
como unidade primária e a lógica alocativa é técnico-clínica. Essa
diferenciação é deliberada: demonstra a versatilidade do pipeline
para acomodar agendas substantivas distintas com o mesmo aparato
técnico subjacente.

O restante do artigo está organizado em cinco seções. A Seção
\ref{sec:teoria} apresenta o marco teórico-clínico. A Seção
\ref{sec:metodo} detalha os aspectos metodológicos. A Seção
\ref{sec:resultados} expõe os resultados empíricos. A Seção
\ref{sec:discussao} discute as implicações. A Seção \ref{sec:conclusao}
apresenta as considerações finais.

% =====================================================================
\section{Marco Teórico-Clínico}\label{sec:teoria}
% =====================================================================

\subsection{Tipologia da incontinência urinária e indicação cirúrgica}

A taxonomia padrão da IU \textemdash{} consolidada na padronização
terminológica IUGA/ICS de 2010 (HAYLEN \textit{et al.}, 2010)
\textemdash{} distingue três tipos principais: \textit{IU de esforço}
(perda associada ao aumento da pressão intra-abdominal por tosse,
espirro, exercício), \textit{IU de urgência} (perda involuntária
acompanhada de urgência miccional) e \textit{IU mista} (combinação
das duas anteriores). A indicação cirúrgica é predominante na IU de
esforço quando há falha do tratamento conservador (fisioterapia
pélvica, comportamental, farmacológico), por se tratar de quadro com
componente anatômico que muitas vezes responde melhor à correção
mecânica.

\subsection{Vias cirúrgicas: abdominal versus vaginal}

A literatura clínica especializada e as diretrizes internacionais
convergem ao priorizar a abordagem vaginal \textemdash{} em particular
os \textit{slings} mediouretrais sintéticos \textemdash{} como
tratamento de primeira linha para IU de esforço em mulheres adultas
(FORD \textit{et al.}, 2017; CHAPPLE \textit{et al.}, 2020). A
abordagem oferece taxa de sucesso comparável à abdominal aberta com
\textbf{menor morbidade}, \textbf{menor permanência hospitalar},
\textbf{menor sangramento intraoperatório} e \textbf{recuperação
funcional mais rápida}. A via abdominal permanece indicada em casos
específicos: cistocele de grande volume, IU recorrente após falha de
\textit{sling}, comorbidades pélvicas associadas ou anatomia
desfavorável à abordagem vaginal.

\subsection{O SUS como provedor cirúrgico de IU}

O SIH-SUS registra, mensalmente, todas as internações hospitalares
faturadas pelo SUS. Cada AIH (Autorização de Internação Hospitalar)
encerra um procedimento principal codificado pela tabela SIGTAP
(Sistema de Gerenciamento da Tabela de Procedimentos, Medicamentos e
OPM) e contém variáveis administrativas (UF, município, hospital,
caráter de atendimento, gestão), clínicas (procedimento, idade, sexo,
óbito hospitalar, dias de permanência) e financeiras (valor total,
serviços hospitalares, serviços profissionais). Essa estrutura permite
análises espaciais e temporais sistemáticas sobre a oferta cirúrgica
do sistema, sob a condição de que o pesquisador seja capaz de
processar os microdados (uma a duas milhões de linhas por UF\,$\times$\,mês).

% =====================================================================
\section{Aspectos Metodológicos}\label{sec:metodo}
% =====================================================================

\subsection{Fontes de dados}

Três fontes oficiais públicas: \textbf{(i)} microdados RD do SIH-SUS
no formato proprietário DBC, disponibilizados via FTP DATASUS em
\url{ftp.datasus.gov.br/dissemin/publicos/SIHSUS/}. Cada arquivo
cobre uma combinação UF\,$\times$\,mês\,$\times$\,ano (ex.:
\texttt{RDSP1503.dbc} = São Paulo, mar/2015). Para a janela
2015\textendash 2020 nas 27 UFs, foram processados aproximadamente
1.944 arquivos. \textbf{(ii)} IBGE/SIDRA Tabela 6579 (população
residente estimada por UF), utilizada para normalizações per capita.
\textbf{(iii)} BCB SGS série 433 (IPCA mensal), utilizada para
deflação a preços de dezembro/2021 \textemdash{} embora o presente
artigo apresente, principalmente, valores nominais por aderência à
tese original.

\subsection{Pipeline de dados}

O pipeline adota a arquitetura medallion (ARMBRUST \textit{et al.},
2021; KUKREJA, 2021) implementada em Apache Spark (ZAHARIA \textit{et
al.}, 2016) sobre Databricks Free Edition com persistência em Delta
Lake. A síntese arquitetural é apresentada na Figura
\ref{fig:architecture}.

\begin{figure}[H]
\centering
\includegraphics[width=0.95\linewidth]{fig02-architecture}
\caption{Pipeline medallion (bronze\,/\,silver\,/\,gold) adotado para o estudo}
\label{fig:architecture}
\par\noindent{\small Fonte: elaboração própria, baseada em
Armbrust \textit{et al.} (2021).}
\end{figure}

\textbf{Bronze} converte cada arquivo DBC em DBF (via PySUS), filtra
linhas pelos SIGTAPs de interesse e materializa Parquet (formato
colunar comprimido). O filtro aplicado já no \textit{convert} reduz
o volume bruto de aproximadamente 150\,GB (todas as colunas, todos
os procedimentos) para alguns megabytes \textemdash{} viabilizando
análise rápida e custo computacional baixo. \textbf{Silver} agrega
por UF\,$\times$\,Ano\,$\times$\,Mês\,$\times$\,Procedimento\,$\times$\,Caráter\,$\times$\,Gestão,
preservando granularidade fina para drill-down. \textbf{Gold}
colapsa para UF\,$\times$\,Ano\,$\times$\,Procedimento e adiciona
deflação por IPCA e per capita por IBGE. Cada camada é versionada
(Delta time travel) e reprodutível.

\subsection{Procedimentos analisados}

O recorte clínico foi limitado aos três SIGTAPs específicos para
tratamento cirúrgico da IU: \texttt{0409010499} (Tratamento Cirúrgico
de Incontinência Urinária Via Abdominal), \texttt{0409070270}
(Tratamento Cirúrgico de Incontinência Urinária Por Via Vaginal) e
\texttt{0409020117} (Tratamento Cirúrgico de Incontinência Urinária
\textemdash{} código genérico, residual no período). A apresentação
dos resultados concentra-se nos dois primeiros, dado que respondem
por aproximadamente 99\% dos registros.

\subsection{Indicadores construídos}

Para cada combinação (UF, Ano, Procedimento) foram construídos:
\textbf{(i)} número de AIH aprovadas; \textbf{(ii)} valor pago total,
hospitalar e profissional, em valores nominais e em \rs 2021;
\textbf{(iii)} permanência média (dias); \textbf{(iv)} taxa de
mortalidade hospitalar; \textbf{(v)} valor médio por AIH; \textbf{(vi)}
AIH por 100\,mil habitantes (per capita ajustado pela escala
apropriada para procedimentos raros).

\subsection{Limitações}

Quatro limitações merecem registro explícito. \textbf{Primeira}, o
SIH-SUS captura apenas a internação faturada \textemdash{} atendimentos
não aprovados, atendimentos no setor privado ou procedimentos
realizados sem AIH não são contemplados. \textbf{Segunda}, a
granularidade espacial está limitada à UF do estabelecimento, não à
residência da paciente; o fluxo interestadual de pacientes não é
diretamente observável neste recorte. \textbf{Terceira}, a taxa de
mortalidade intra-hospitalar foi excluída da análise central pela
alta esparsidade dos dados (a maioria das células UF\,$\times$\,Ano
registra zero óbitos). \textbf{Quarta}, a análise é descritiva e não
pretende identificar relações causais entre as variáveis institucionais
e os volumes observados.

% =====================================================================
\section{Resultados}\label{sec:resultados}
% =====================================================================

\subsection{Volumes nacionais e razão entre vias}

Nos seis anos analisados (2015\textendash 2020), o SUS aprovou
\textbf{4.472} AIH para tratamento cirúrgico de IU via abdominal e
\textbf{30.688} AIH para a via vaginal. A razão entre vias é de
aproximadamente \textbf{6{,}86:1} em favor da via vaginal
\textemdash{} contraste apresentado na Figura
\ref{fig:volume-contrast}.

\begin{figure}[H]
\centering
\includegraphics[width=0.95\linewidth]{fig05-volume-contrast}
\caption{Contraste de volume total entre as duas vias cirúrgicas, AIH acumuladas 2015\textendash 2020}
\label{fig:volume-contrast}
\par\noindent{\small Fonte: elaboração própria, microdados SIH-SUS/DATASUS.}
\end{figure}

A evolução anual desagregada é apresentada na Figura
\ref{fig:evolution-aih}. Três regularidades chamam a atenção.
\textbf{Primeira}, o volume da via vaginal é estruturalmente superior
em todos os anos da série, com diferencial mantido em torno de
7\textendash 8 vezes em condições pré-pandemia. \textbf{Segunda}, a
via abdominal apresenta tendência declinante já a partir de 2016
(de 1.020 para 759 AIH/ano até 2019), possivelmente refletindo a
consolidação de protocolos clínicos que privilegiam a abordagem
vaginal. \textbf{Terceira}, em 2020 ambas as vias sofrem queda da
ordem de 50\% (ABD: 759\,$\rightarrow$\,384; VAG: 5.959\,$\rightarrow$\,2.566),
inflexão consistente com o adiamento de cirurgias eletivas durante a
pandemia de COVID-19 \textemdash{} fenômeno amplamente documentado
no SUS (NUNES \textit{et al.}, 2022).

\begin{figure}[H]
\centering
\includegraphics[width=0.95\linewidth]{fig03-evolution-aih}
\caption{Volume anual de AIH aprovadas, 2015\textendash 2020, por via cirúrgica}
\label{fig:evolution-aih}
\par\noindent{\small Fonte: elaboração própria, microdados SIH-SUS/DATASUS.}
\end{figure}

\subsection{Despesa pública: evolução e composição}

O gasto público total acumulado em 2015\textendash 2020, em valores
nominais (\rs correntes), atingiu \textbf{\rs 1.979.143,90} para a
via abdominal e \textbf{\rs 13.035.156,65} para a via vaginal,
totalizando aproximadamente \textbf{\rs 15{,}0 milhões}. A
distribuição anual é apresentada na Figura \ref{fig:evolution-valor}.
Em termos relativos ao orçamento global do SUS, o montante é modesto;
em termos absolutos, contudo, é significativo dado o caráter
altamente eletivo dos procedimentos.

\begin{figure}[H]
\centering
\includegraphics[width=0.95\linewidth]{fig04-evolution-valor}
\caption{Valor pago anual (\rs nominais, milhares) por via cirúrgica}
\label{fig:evolution-valor}
\par\noindent{\small Fonte: elaboração própria. Valores nominais \textemdash{}
para análise em \rs 2021 deflacionados, ver tabela gold do pipeline.}
\end{figure}

A Tabela \ref{tab:evolucao-uropro} consolida a série de AIH e valores
para o quinquênio.

\begin{table}[H]
\centering
\caption{AIH aprovadas e valor pago por via cirúrgica, Brasil, 2015\textendash 2020}
\begin{tabular}{rrrrr}
\toprule
\textbf{Ano} & \textbf{AIH ABD} & \textbf{AIH VAG} & \textbf{\rs ABD (mil)} & \textbf{\rs VAG (mil)}\\
\midrule
2015 & 1.020 &  5.722 & 475{,}3 & 2.586{,}6 \\
2016 &   789 &  5.467 & 328{,}8 & 2.188{,}4 \\
2017 &   760 &  5.419 & 326{,}5 & 2.228{,}2 \\
2018 &   760 &  5.555 & 353{,}3 & 2.472{,}2 \\
2019 &   759 &  5.959 & 338{,}5 & 2.566{,}2 \\
2020 &   384 &  2.566 & 156{,}8 &   993{,}5 \\
\midrule
\textbf{Total} & \textbf{4.472} & \textbf{30.688} & \textbf{1.979{,}1} & \textbf{13.035{,}2} \\
\bottomrule
\end{tabular}
\par\noindent{\small Fonte: elaboração própria, microdados SIH-SUS via TabNet (Tatieli, 2022).
Valores em \rs nominais correntes do ano de processamento.}
\label{tab:evolucao-uropro}
\end{table}

\subsection{Distribuição geográfica}

A Figura \ref{fig:top10-vag} apresenta as dez UFs com maior volume
acumulado de AIH para a via vaginal e a Figura \ref{fig:top10-abd}
faz o mesmo para a via abdominal. \textbf{São Paulo lidera em ambos
os recortes}, mas com magnitudes muito distintas: 29{,}8\% do total
nacional na via abdominal (1.335\,/\,4.472) e \textbf{70{,}3\%} na
via vaginal (21.588\,/\,30.688). A leitura mais imediata é a de
forte concentração assistencial em hospitais-referência paulistas
(HC-FMUSP, HSP-EPM, ICESP).

\begin{figure}[H]
\centering
\includegraphics[width=0.85\linewidth]{fig07-top10-vag}
\caption{Top-10 UFs em AIH cirúrgicas via vaginal (acumulado 2015\textendash 2020)}
\label{fig:top10-vag}
\par\noindent{\small Fonte: elaboração própria, microdados SIH-SUS/DATASUS.}
\end{figure}

\begin{figure}[H]
\centering
\includegraphics[width=0.85\linewidth]{fig06-top10-abd}
\caption{Top-10 UFs em AIH cirúrgicas via abdominal (acumulado 2015\textendash 2020)}
\label{fig:top10-abd}
\par\noindent{\small Fonte: elaboração própria, microdados SIH-SUS/DATASUS.}
\end{figure}

A visualização integral por UF, em mapa coroplético, é apresentada na
Figura \ref{fig:choro-vag} (via vaginal) e na Figura \ref{fig:choro-abd}
(via abdominal). A hegemonia paulista é nítida, especialmente para a
via vaginal. Estados como \textbf{Alagoas} (sexto colocado em ABD com
269 AIH, mas vigésimo-segundo em VAG) e \textbf{Paraná} (oitavo em ABD,
mas segundo em VAG) sugerem padrões heterogêneos de adoção tecnológica
e organização de serviços que merecem investigação qualitativa
complementar.

\begin{figure}[H]
\centering
\includegraphics[width=0.78\linewidth]{fig08-choropleth-vag}
\caption{Mapa coroplético \textemdash{} AIH cirúrgicas via vaginal por UF (acumulado 2015\textendash 2020)}
\label{fig:choro-vag}
\par\noindent{\small Fonte: elaboração própria. GeoJSON IBGE + microdados SIH-SUS.}
\end{figure}

\begin{figure}[H]
\centering
\includegraphics[width=0.78\linewidth]{fig09-choropleth-abd}
\caption{Mapa coroplético \textemdash{} AIH cirúrgicas via abdominal por UF (acumulado 2015\textendash 2020)}
\label{fig:choro-abd}
\par\noindent{\small Fonte: elaboração própria. GeoJSON IBGE + microdados SIH-SUS.}
\end{figure}

A análise per capita (Figura \ref{fig:percap-vag}) atenua \textemdash{}
sem eliminar \textemdash{} a concentração observada em valores
absolutos. Estados de menor população, como Amapá e Roraima, sobem
proporcionalmente no ranking quando normalizados por habitante, ainda
que sigam abaixo da média nacional em valores absolutos.

\begin{figure}[H]
\centering
\includegraphics[width=0.7\linewidth]{fig12-per-capita-vag}
\caption{Ranking horizontal: AIH via vaginal por 100 mil habitantes (acumulado 2015\textendash 2020)}
\label{fig:percap-vag}
\par\noindent{\small Fonte: elaboração própria, SIH-SUS + IBGE/SIDRA Tabela 6579 (pop.\ 2020).}
\end{figure}

A interpretação geográfica deve considerar que o SIH-SUS registra a
AIH no estabelecimento que realizou o procedimento, não na residência
da paciente: a concentração em SP reflete parcialmente o fluxo
interestadual de pacientes encaminhadas a hospitais-referência. Para
fins de planejamento de saúde pública, ambas as leituras são
relevantes: a leitura absoluta dimensiona a oferta efetiva regional,
e a leitura per capita dimensiona o acesso da população local.

\subsection{Permanência hospitalar}

A permanência hospitalar média é diferente entre as vias e consistente
com as expectativas clínicas (Figura \ref{fig:permanencia}). Para o
conjunto do Brasil no período 2015\textendash 2020, a via abdominal
apresenta média de \textbf{3{,}47 dias} de internação contra
\textbf{2{,}16 dias} da via vaginal \textemdash{} diferença de
\textbf{37{,}8\%} a favor desta última. Esse resultado é compatível
com a literatura clínica que reporta recuperação mais rápida e menores
complicações pós-operatórias para a abordagem vaginal (FORD
\textit{et al.}, 2017).

\begin{figure}[H]
\centering
\includegraphics[width=0.85\linewidth]{fig10-permanencia}
\caption{Permanência hospitalar média (dias) por via, média 2015\textendash 2020}
\label{fig:permanencia}
\par\noindent{\small Fonte: elaboração própria, microdados SIH-SUS/DATASUS.}
\end{figure}

\subsection{Custo médio por AIH}

Apesar da diferença de permanência, o custo médio por AIH é
\textbf{praticamente equivalente} entre as duas vias (Figura
\ref{fig:custo}): \textbf{\rs 442{,}61} (ABD) vs.\ \textbf{\rs 424{,}79}
(VAG), diferença de apenas \textbf{+4{,}2\%} em favor da via abdominal.
A interpretação mais plausível é que a remuneração SIGTAP-SUS opera em
base fixa por procedimento, com permanência subjacente refletida apenas
parcialmente nos valores hospitalares (\texttt{VAL\_SH}). Sob essa
lógica, a via vaginal é \textbf{significativamente mais eficiente do
ponto de vista do custo social}: menor permanência, menor uso de
recursos hospitalares, recuperação mais rápida da paciente, com
remuneração praticamente idêntica.

\begin{figure}[H]
\centering
\includegraphics[width=0.85\linewidth]{fig11-custo-medio}
\caption{Custo médio por AIH por via, valores nominais 2015\textendash 2020}
\label{fig:custo}
\par\noindent{\small Fonte: elaboração própria, microdados SIH-SUS/DATASUS.}
\end{figure}

\subsection{Efeito da pandemia (2020)}

O ano de 2020 apresenta queda abrupta de aproximadamente 50\% em ambas
as vias (Figura \ref{fig:evolution-aih}), sendo a abdominal um pouco
menos afetada (\textendash 49{,}4\%, de 759 em 2019 para 384 em 2020)
que a vaginal (\textendash 56{,}9\%, de 5.959 para 2.566). Esse padrão
é consistente com a evidência disponível sobre o impacto da pandemia
sobre cirurgias eletivas no SUS (NUNES \textit{et al.}, 2022): a
suspensão temporária e a realocação de leitos para pacientes com
COVID-19 levaram ao adiamento generalizado de procedimentos
não-emergenciais, situação especialmente acentuada em hospitais de
alta complexidade \textemdash{} justamente os que concentram a oferta
para tratamento de IU.

% =====================================================================
\section{Discussão}\label{sec:discussao}
% =====================================================================

\subsection{Aderência às boas práticas clínicas}

O perfil agregado nacional \textemdash{} predomínio robusto e estável
da via vaginal, com proporção próxima a 7:1 \textemdash{} é compatível
com as recomendações internacionais de uroginecologia (FORD
\textit{et al.}, 2017; CHAPPLE \textit{et al.}, 2020). O achado, em
si, é uma contribuição positiva da pesquisa: documenta empiricamente
uma realidade frequentemente desconhecida tanto pelo público leigo
quanto pelo gestor hospitalar de outras especialidades, contribuindo
para a superação do estigma associado ao tratamento da IU.

\subsection{Heterogeneidade regional e implicações para a Enfermagem}

A hegemonia paulista é o achado mais visível, mas a heterogeneidade
regional revela-se mais informativa quando dissecada por via. Estados
como Alagoas (sexto na via abdominal, mas vigésimo-segundo na vaginal)
ou Paraná (oitavo na abdominal, segundo na vaginal) sugerem padrões de
adoção tecnológica e organização de serviços que merecem investigação
qualitativa complementar. Para a Enfermagem uroginecológica e a
gestão hospitalar, o mapeamento empírico dessas heterogeneidades
fornece base objetiva para a discussão sobre alocação de equipes,
capacitação técnica e dimensionamento da rede de referência.

\subsection{Custo, permanência e remuneração: implicações para a gestão}

O paradoxo do custo médio similar com permanência muito distinta
sugere ineficiência alocativa potencial na tabela SIGTAP-SUS:
remunerar igualmente procedimentos com perfis de uso de recursos tão
distintos cria incentivo perverso à adoção da abordagem mais
custosa para o sistema (maior permanência, maior uso de leito, maior
risco de complicações pós-operatórias). Uma agenda de política
pública relevante seria a revisão da tabela SIGTAP para
\textit{slings} mediouretrais com remuneração diferencial, capaz de
sinalizar mais explicitamente a preferência pela abordagem
minimamente invasiva.

\subsection{Limitações da inferência}

Quatro limitações merecem ser reiteradas. \textbf{Primeira}, os
resultados são \textbf{descritivos}: não permitem identificação causal
isolada do efeito de cada variável institucional. \textbf{Segunda}, o
SIH captura apenas a internação faturada \textemdash{} não permite
afirmar nada sobre a IU não tratada cirurgicamente, atendimentos no
setor privado ou procedimentos sem AIH. \textbf{Terceira}, a
agregação ao nível UF do estabelecimento mascara fluxos
interestaduais de pacientes que precisariam de cruzamento com
\texttt{MUNIC\_RES} para serem desambiguados. \textbf{Quarta}, a
janela 2015\textendash 2020 é a janela da pesquisa original e a
mais recente disponível com cobertura completa nos arquivos TabNet
quando a tese foi defendida; a integração ao pipeline aberto permite
agora estender a análise para 2021\textendash 2025 conforme o
\textit{backlog} pós-pandemia for absorvido.

% =====================================================================
\section{Considerações Finais}\label{sec:conclusao}
% =====================================================================

\subsection{O que descobrimos}

\textbf{(i) Predomínio robusto e estável da via vaginal.} A razão entre
AIH vaginais e abdominais é de aproximadamente \textbf{6{,}86:1} em
todo o período (30.688 vs.\ 4.472 AIH em 2015\textendash 2020), padrão
mantido em todos os anos da série e compatível com as recomendações
internacionais de uroginecologia. O achado oferece evidência empírica
de aderência do SUS, em seus serviços-referência, às boas práticas
clínicas contemporâneas.

\textbf{(ii) Permanência hospitalar 37{,}8\% menor na via vaginal.} A
média do Brasil em 2015\textendash 2020 é de 3{,}47 dias para a via
abdominal contra 2{,}16 dias para a vaginal, consistente com a menor
invasividade da abordagem vaginal e seu perfil de recuperação mais
rápida. Para a Enfermagem hospitalar, esse diferencial tem implicações
diretas sobre dimensionamento de equipe, gestão de leitos e
planejamento de alta.

\textbf{(iii) Custo médio por AIH praticamente idêntico entre as
vias.} \rs 442{,}61 (ABD) vs.\ \rs 424{,}79 (VAG) \textemdash{}
diferença de apenas +4{,}2\%, apesar da permanência hospitalar 60\%
mais longa na ABD. Em termos de \textit{custo social}, a via vaginal
entrega o mesmo desfecho remunerado a custo menor.

\textbf{(iv) Concentração geográfica extrema, especialmente na via
vaginal.} São Paulo concentra 29{,}8\% das AIH abdominais e
\textbf{70{,}3\%} das vaginais, refletindo provável centralização em
hospitais-referência. Para a gestão pública de saúde, esse padrão
sinaliza tanto fluxo interestadual de pacientes quanto potencial
subdimensionamento da oferta em estados de média população.

\textbf{(v) Queda de \,$\sim$\,50\% em 2020 atribuível à pandemia.} O
efeito da pandemia de COVID-19 sobre cirurgias eletivas é documentado
de forma quantitativa: ABD \textendash 49{,}4\%, VAG \textendash 56{,}9\%
entre 2019 e 2020. A magnitude do \textit{backlog} resultante e sua
duração permanecem em aberto \textemdash{} investigação que se torna
possível com o pipeline atualizável aqui apresentado.

\subsection{Agenda de pesquisa}

Quatro frentes destacam-se como prioritárias: \textbf{(a)} extensão da
janela analítica para 2021\textendash 2025 com foco no \textit{backlog}
pós-pandemia; \textbf{(b)} cruzamento com \texttt{MUNIC\_RES} para
desambiguar concentração assistencial vs.\ fluxo interestadual;
\textbf{(c)} integração com a vertical CNES já disponível no
\textit{Mirante dos Dados} (CHALHOUB, 2026c, Working Paper futuro),
permitindo testar associações entre disponibilidade de equipamentos
uroginecológicos e volumes cirúrgicos por UF; e \textbf{(d)} avaliação
prospectiva de desfechos clínicos (re-internação, mortalidade tardia)
por via, recurso que demanda articulação com bases SIM/SINASC
(mortalidade) hoje fora do escopo do Mirante.

\subsection{Contribuição metodológica}

Para além dos achados substantivos, a contribuição principal é de
\textbf{infraestrutura}: o pipeline aqui apresentado é integralmente
\textit{open-source}, segue os princípios FAIR de gestão de dados
científicos (WILKINSON \textit{et al.}, 2016), é versionado em Git
com \textit{refresh} mensal automatizado, e é extensível a qualquer
procedimento cirúrgico do SIGTAP-SUS por modificação de um único
parâmetro (\texttt{procs\_filter}). Em conjunto com os Working
Papers n.\,1 (CHALHOUB, 2026a) e n.\,2 (CHALHOUB, 2026b), o presente
trabalho demonstra a viabilidade de produzir múltiplas pesquisas
substantivas a partir de uma única base consolidada \textemdash{}
reduzindo o custo marginal de investigação para profissionais de
Enfermagem, gestores hospitalares, jornalistas e pesquisadores em
saúde coletiva.

% =====================================================================
\section*{Referências}\label{sec:ref}
\addcontentsline{toc}{section}{Referências}
% =====================================================================
\begin{singlespace}\small
\noindent\hangindent=1.25cm\hangafter=1
ARMBRUST, M.; GHODSI, A.; XIN, R.; ZAHARIA, M. Lakehouse: a new
generation of open platforms that unify data warehousing and advanced
analytics. In: \textbf{11th CIDR}, 2021.\par\smallskip

\noindent\hangindent=1.25cm\hangafter=1
BANCO CENTRAL DO BRASIL (BCB). \textbf{Sistema Gerenciador de Séries
Temporais}, série 433: IPCA. Disponível em:
\url{https://api.bcb.gov.br/dados/serie/bcdata.sgs.433/dados}.
Acesso em: abr. 2026.\par\smallskip

\noindent\hangindent=1.25cm\hangafter=1
BRASIL. Ministério da Saúde. \textbf{Sistema de Informações
Hospitalares do SUS (SIH-SUS)}. Microdados. Brasília: DATASUS.
Disponível em:
\url{https://datasus.saude.gov.br/transferencia-de-arquivos/}.
Acesso em: abr. 2026.\par\smallskip

\noindent\hangindent=1.25cm\hangafter=1
BRASIL. Ministério da Saúde. \textbf{Tabela de Procedimentos,
Medicamentos e OPM do SUS \textemdash{} SIGTAP}. Brasília: DATASUS,
2024.\par\smallskip

\noindent\hangindent=1.25cm\hangafter=1
CHALHOUB, L. Emendas Parlamentares no Orçamento Federal Brasileiro
(2015\textendash 2025): distribuição espacial, execução orçamentária
e efeitos das mudanças institucionais recentes.
\textbf{Mirante dos Dados \textemdash{} Working Paper n. 1}, abr.
2026a.\par\smallskip

\noindent\hangindent=1.25cm\hangafter=1
CHALHOUB, L. Programa Bolsa Família, Auxílio Brasil e Novo Bolsa
Família (2013\textendash 2025): transformações institucionais,
expansão da cobertura e desigualdade territorial.
\textbf{Mirante dos Dados \textemdash{} Working Paper n. 2}, abr.
2026b.\par\smallskip

\noindent\hangindent=1.25cm\hangafter=1
CHAPPLE, C. R.; CRUZ, F.; DESCHAMPS, C.; HAYLEN, B. T. \textit{et al.}
Consensus statement of the European Association of Urology on the
surgical treatment of urinary incontinence in women.
\textbf{European Urology}, v. 78, n. 5, p. 643\textendash 656, 2020.\par\smallskip

\noindent\hangindent=1.25cm\hangafter=1
FORD, A. A.; ROGERSON, L.; CODY, J. D.; OGAH, J. Mid-urethral sling
operations for stress urinary incontinence in women.
\textbf{Cochrane Database of Systematic Reviews}, n. 7, CD006375,
2017.\par\smallskip

\noindent\hangindent=1.25cm\hangafter=1
HAYLEN, B. T.; DE RIDDER, D.; FREEMAN, R. M. \textit{et al.} An
International Urogynecological Association
(IUGA)\,/\,International Continence Society (ICS) joint report on
the terminology for female pelvic floor dysfunction.
\textbf{Neurourology and Urodynamics}, v. 29, n. 1,
p. 4\textendash 20, 2010.\par\smallskip

\noindent\hangindent=1.25cm\hangafter=1
HU, T.-W.; WAGNER, T. H.; BENTKOVER, J. D. \textit{et al.} Costs of
urinary incontinence and overactive bladder in the United States: a
comparative study. \textbf{Urology}, v. 63, n. 3,
p. 461\textendash 465, 2003.\par\smallskip

\noindent\hangindent=1.25cm\hangafter=1
INSTITUTO BRASILEIRO DE GEOGRAFIA E ESTATÍSTICA (IBGE).
\textbf{SIDRA \textemdash{} Sistema IBGE de Recuperação Automática}.
Tabela 6579: Estimativas anuais da população residente. Disponível em:
\url{https://sidra.ibge.gov.br/tabela/6579}. Acesso em: abr. 2026.
\par\smallskip

\noindent\hangindent=1.25cm\hangafter=1
KUKREJA, M. \textbf{Data Engineering with Apache Spark, Delta Lake,
and Lakehouse}. Packt Publishing, 2021.\par\smallskip

\noindent\hangindent=1.25cm\hangafter=1
NUNES, B. P.; FLORES, T. R.; MIELKE, G. I. \textit{et al.} The
COVID-19 pandemic and elective surgeries in the Brazilian Unified
Health System (SUS). \textbf{Public Health}, v. 211,
p. 109\textendash 115, 2022.\par\smallskip

\noindent\hangindent=1.25cm\hangafter=1
TATIELI DA SILVA. \textbf{Análise Descritiva da Cirurgia para
Tratamento da Incontinência Urinária no SUS, 2015\textendash 2020}.
Trabalho de Conclusão de Curso (Especialização em Enfermagem),
2022.\par\smallskip

\noindent\hangindent=1.25cm\hangafter=1
WILKINSON, M. D. \textit{et al.} The FAIR Guiding Principles for
scientific data management and stewardship. \textbf{Scientific Data},
v. 3, 2016. DOI: 10.1038/sdata.2016.18.\par\smallskip

\noindent\hangindent=1.25cm\hangafter=1
ZAHARIA, M. \textit{et al.} Apache Spark: a unified engine for big
data processing. \textbf{Communications of the ACM}, v. 59, n. 11,
2016.\par\smallskip
\end{singlespace}

\vspace{1cm}

\noindent\textbf{Citar como:}\\
DA SILVA, T.; CHALHOUB, L. Tratamento Cirúrgico da Incontinência
Urinária no Sistema Único de Saúde (2015\textendash 2020): volumes,
despesa, permanência e distribuição geográfica por via de acesso.
\textit{Mirante dos Dados \textemdash{} Working Paper n. 3}, v. 1.0,
abr. 2026. Disponível em:
\url{https://leonardochalhoub.github.io/mirante-dos-dados-br/}.

\vspace{0.5cm}

\noindent\textbf{Reconhecimento:} a análise empírica original deste
artigo foi conduzida por Tatieli da Silva como trabalho de conclusão
de especialização em Enfermagem (2022). A presente versão reproduz os
resultados a partir dos microdados SIH-RD, adicionando dimensões de
análise (custo médio, COVID-19, comparação direta entre vias) e
integrando o trabalho ao pipeline aberto Mirante dos Dados.

\vspace{0.5cm}

\noindent\textbf{Licença:} dados consolidados (camada Gold) e código
do pipeline distribuídos sob licença MIT. Texto deste artigo
distribuído sob licença Creative Commons Atribuição 4.0 Internacional
(CC BY 4.0).

\end{document}
